\lesson{5}{Dec 01 2021 Wed (19:09:04)}{Polynomial Identities and Proofs}{Unit 3}

\begin{definition}[Algebraic Proofs]
    \begin{itemize}
        \item \bf{Polynomial Identities} can be proven to be true by simplifying the identity through application of \bf{Algebraic Theorems} and \bf{Principles}.
        \item Start with the side of the identity that can be simplified the easiest.
        \item Sometimes, following a \it{“clue”} will lead to a dead-end in your \bf{Proof}. Do not give up. Just follow a different \it{“clue”}. The more practice you have with proofs, the more you will be able to predict the dead-ends.
    \end{itemize}
\end{definition}

\begin{example}

\end{example}

\subsubsection*{Application to Numerical Relationships}

Polynomial identities apply to more than just polynomials. Replacing the variables with numbers can help prove numerical relationships as well.

\newpage
